\chapter{总结与展望}
\label{chap:conclusion}

\section{工作总结}

本文围绕 \textbf{Survivor And Fight}，在 LayaAir 3.x 与 TypeScript 上完成 Roguelite 生存战斗的分析、设计、实现与测试，对应关系如下：

\begin{enumerate}[label=(\arabic*)]
  \item \textbf{需求与选型}（第~\ref{chap:intro}--\ref{chap:related}~章）：梳理引擎架构、ECS、人群仿真与 H5 性能文献；确定 LayaAir、轻量 ECS、MVC UI、配表战斗的路线。
  \item \textbf{总体设计}（第~\ref{chap:architecture}~章）：给出分层结构、ECS/MVC/Worker 方案，以及时序、可靠性、容错和扩展方式。
  \item \textbf{实现}（第~\ref{chap:modules}~章）：以 \texttt{SimpleEcsDemo} 为根，实现技能链、子弹、怪物 AI、跑图、Tab 装配和主循环。
  \item \textbf{性能与测试}（第~\ref{chap:performance}--\ref{chap:test}~章）：千怪场景池化明显压低 P95；图集抬高重度弹幕波 FPS；用例和脚本验证核心契约。
\end{enumerate}

写作中参考 SyDRA\cite{ullmann2025sydra}、Gregory 分层\cite{gregory2018game} 和 Unity 人群实践\cite{cantrell2018unity}，把模块化、重算分离和 V\&V 用到 H5 课设原型上。系统能演示菜单 $\rightarrow$ 跑图/选关 $\rightarrow$ 战斗 $\rightarrow$ 升级装配 $\rightarrow$ 死亡重启，达到毕设预期。

\section{不足与展望}

\begin{enumerate}[label=\arabic*.]
  \item \textbf{内容与表现}：法杖三槽 UI、部分 Effect 特效和音效仍缺；
  \item \textbf{架构}：ECS 可改 Archetype/Chunk 提同屏上限\cite{ullmann2023coupling}；碰撞可加空间哈希；
  \item \textbf{测试}：补 Worker「池+Worker、已预热」同窗消融和 UI 自动化\cite{petty2010vv}；
  \item \textbf{产品化}：局外存档、元进度、帧率与同屏怪上限设置；
  \item \textbf{发布}：补美术与适龄说明后静态托管。
\end{enumerate}

\section{社会、健康与文化影响}

战斗内容为虚构，怪物和血条偏抽象，避免写实血腥。原型无付费、无诱导分享、不采敏感数据。若上网发布，需遵守网游和未成年人保护规定，标适龄和时长；长时间玩注意眼疲劳，Tab 暂停和跑图休息节点用来调节节奏\cite{petty2010vv}。

\section{绿色节能与兼容性}

池化、暂停和 Worker 降低平均 CPU，对续航和发热有帮助\cite{gregory2018game}。配表双通道和 Worker 回退提高浏览器兼容性；2D 对 GPU 压力小于大型 3D。后续可在设置里加帧率上限和同屏怪上限，做低功耗档。

\section{本章小结}

Survivor And Fight 说明在 H5 上用轻量 ECS、配表和可关断的性能手段做 Roguelite 原型是可行的。论文叙述与仓库、测试表一致；后续可补 Worker 定量、自动化测试和表现层，从答辩 Demo 往可发布小游戏再走一步。
